% !TEX root = z-main.tex

El avance de las tecnologías digitales transformó las formas de preservar, estudiar y difundir el patrimonio cultural. En este contexto, el modelado tridimensional (\gls{3d}), junto con la realidad aumentada (\gls{RA}), se consolidaron como herramientas fundamentales en la arqueología y la museología. Estas tecnologías permitieron recrear digitalmente sitios históricos, incluidos aquellos parcialmente desaparecidos o, como Kaminaljuyú, aún enterrados.

Kaminaljuyú, ubicado en el altiplano central guatemalteco, fue una de las ciudades más importantes del período preclásico mesoamericano. Su ocupación abarcó aproximadamente desde el año 1200 a. C. hasta el 900 d. C., extendiéndose sobre una superficie de alrededor de cinco kilómetros cuadrados y conformándose por más de 230 montículos. Debido al crecimiento urbano de la Ciudad de Guatemala, gran parte del sitio fue destruido o cubierto por construcciones modernas, conservándose únicamente cerca de cuarenta montículos, de los cuales nueve permanecen resguardados dentro del Parque Arqueológico Kaminaljuyú. Entre las estructuras más representativas que aún se mantienen visibles se encuentran la Acrópolis (C-II-4) y la Estructura E de La Palangana. Asimismo, hay estructuras que se encuentran cubiertas por la vegetación, dificultando su interpretación volumétrica. Esta pérdida progresiva del paisaje arqueológico ha planteado la necesidad de recurrir a herramientas digitales para documentar, preservar y facilitar la comprensión visual de lo que permanece del sitio.

En este marco, el trabajo tuvo como propósito modelar en \gls{3d} una selección de estructuras de Kaminaljuyú, con el fin de integrarlas en una aplicación de realidad aumentada que posibilitara su visualización e interacción digital. El desarrollo se basó en el uso del software Blender, empleando datos técnicos proporcionados por el área de arqueología, la cual validó cada modelo. La planificación del proyecto se gestionó de manera centralizada a través de la plataforma Notion, donde se documentaron los avances y se asignaron las tareas correspondientes.

El proceso incluyó la recopilación y análisis de información arqueológica, la construcción progresiva de los modelos desde las bases hasta las secciones superiores, la incorporación de texturas seleccionadas desde Poly Haven y la validación arqueológica respectiva. Los modelos finales fueron almacenados en un repositorio privado en GitHub en formatos compatibles, mientras que los archivos fuente se conservaron en OneDrive para garantizar su preservación.

La investigación demostró que el uso de herramientas digitales contribuyó a la preservación y difusión del patrimonio maya guatemalteco. Asimismo, se planteó que los resultados no se limitaron al ámbito académico o técnico, sino que también ofrecieron beneficios para instituciones culturales, turísticas y educativas interesadas en promover el conocimiento del pasado mediante experiencias tecnológicas inmersivas.

El trabajo se estructura en capítulos que abordan el contexto arqueológico, la metodología de modelado y desarrollo, los resultados obtenidos y la discusión final.

En conjunto, el trabajo evidenció que la integración de tecnologías tridimensionales y de realidad aumentada constituye una herramienta viable para la conservación digital del patrimonio arqueológico guatemalteco.